\chapter{What you will have at the end}
\label{ch:before}

By the end of this guide you will have a personal website at an address like \texttt{https://your-name.github.io}. It will have a home page, a page for your work and a CV page. It will work on phones and offer light and dark themes. It costs nothing to run, needs no server, and is yours to change.

\begin{figure}[H]\centering\includegraphics[width=\linewidth]{web_template_live.png}\caption{The site you will start from, before you change the words. You can open it at \texttt{ahmadammar.github.io/academic-site-template}.}\end{figure}

\section{Who this is for}

You need no experience with code. If you can install a program and follow steps in order, you can do this. Set aside about an hour for the first time, most of it spent on installing things. After that, changing a page takes a few minutes.

Everything you need is in this guide, including in the printed copy. Web addresses are written in full so you can type them.

\section{How it works, in one paragraph}

GitHub offers a free service called GitHub Pages. It takes the files in a repository, which is a project folder stored on GitHub, and shows them as a website. So the job has four parts: get a copy of the template, change the words on your own computer, send the changes to GitHub, and switch Pages on.

\input{shared_git_words}

\input{shared_terminal}
